\documentclass[11pt]{article}
\usepackage[margin=25mm]{geometry}
\usepackage{fontspec}
\setmainfont{Latin Modern Roman}
\usepackage{amsmath,amssymb,mathtools}
\usepackage{unicode-math}
\setmathfont{Latin Modern Math}
\usepackage{xurl,hyperref}
\hypersetup{colorlinks=true,urlcolor=blue}
\setcounter{secnumdepth}{0}
\setlength{\parindent}{0pt}
\setlength{\parskip}{5pt}
\setlength{\emergencystretch}{3em}
\newcommand{\sref}[2]{\hyperlink{source:#1:#2}{[S#2]}}
\newcommand{\cataloguescope}[1]{\paragraph{Recorded target.}#1}
\newcommand{\R}{\mathbb R}
\newcommand{\C}{\mathbb C}
\newcommand{\N}{\mathbb N}
\newcommand{\Z}{\mathbb Z}
\newcommand{\Q}{\mathbb Q}
\newcommand{\F}{\mathbb F}
\newcommand{\eps}{\varepsilon}
\begin{document}
% UnsolvedMath ID 20002542; source catalogue code AIM-PHYSICS-0050
\setcounter{section}{20002541}
\section{Effective full-isotropy constraints for time-space-label symmetries}\label{20002542}
{\small UnsolvedMath ID 20002542 · Mathematical Physics}
\subsection{Definitions and mathematical statement}

\paragraph{Shared notation.}$\mathbb N=\{1,2,\ldots\}$; the standard number fields have their usual meanings. All additional conventions are stated locally.


\paragraph{Definitions and notation.}In the Newtonian $n$-body problem, positions $q_i\in\mathbb R^d$ and masses $m_i>0$ obey $m_i\ddot q_i=\nabla_{q_i}U$, where $U(q)=\sum_{i<j}m_im_j/|q_i-q_j|$ and the gravitational constant is normalized to one. Here $T=\tfrac12\sum_i m_i|\dot q_i|^2$, $E=T-U$, and a collision has $q_i=q_j$. The center of mass can be fixed at zero; moment of inertia is $I=\sum_i m_i|q_i|^2$. For periodic motion in $\mathbb R^d$, time reparametrization by a rotation or reflection of the circle gives an $O(2)$ action; spatial orthogonal transformations give $O(d)$, and label permutations give $S_n$. The stabilizer consists of triples preserving the loop. With unequal masses, a label symmetry must respect the mass assignment to preserve the equations.

\paragraph{Catalogue wording retained for traceability.}
Which subgroups of $O(2)\times O(d)\times S_n$ are realized as symmetries of solutions? Can the symmetry group arise in different ways, as the second question in the source asks?
\paragraph{Scope and source qualification.}Realization as a contained symmetry group and realization as the full stabilizer are different; the original question does not explicitly choose one. The classification and multiplicity-of-realization request must both be addressed without inventing that choice. The catalogue title/status was paired with a machine-generated research attempt; that attempt is not peer-reviewed evidence of settlement. The statement covers the original source questions rather than only the title’s subcase.
\subsection{Short English statement}
Determine the time, space and label symmetry groups realizable by n-body solutions.
\subsection{Research attempt}
\paragraph{Model, date and investigated scope.}
OpenAI GPT-6 Astra (configured model \texttt{gpt-6-astra}), ultra reasoning; generated 11 October 2026. Proved the collision-free pure-label kernel obstruction and displayed continuous extra isotropy of regular polygon choreographies. No classification or multiplicity-of-realization theorem.

\paragraph{Sources and literature check.}
Read the supplied complete definition, catalogue wording, scope qualifications and source list for this canonical ID. The original-source citations are inherited as supplied references; no fresh primary-source or novelty search was performed for this attempt. No current open-status claim is inferred from this calculation.

\paragraph{An effective-kernel constraint.}For a collision-free periodic solution, suppose a symmetry triple acts trivially in time and space and only permutes labels by \(\sigma\). Invariance requires \(q_{\sigma(i)}(t)=q_i(t)\) for every i,t. Collision freedom forces \(\sigma(i)=i\) for all i. Thus the full stabilizer has trivial intersection with the pure-label subgroup. In particular an abstract proposed symmetry group with a nontrivial pure-label kernel cannot be realized by a collision-free labeled solution.
\paragraph{Continuous isotropy of a circle choreography.}For the equal-mass rotating regular n-gon, with \(q_j(t)=R e^{i(\omega t+2\pi j/n)}\), every time shift \(t\mapsto t+s\) can be canceled by spatial rotation \(e^{-i\omega s}\). These triples form a continuous circle subgroup of the stabilizer with identity label permutation. The discrete time shift \(P/n\) can also be matched with a cyclic label permutation. Hence specifying only a finite cyclic symmetry subgroup does not identify the full stabilizer; a realization may carry additional continuous symmetries.
\paragraph{Mass and scope checks.}Permutations used as symmetries of the equations must preserve the assigned masses. Spatial transformations acting trivially on a lower-dimensional trace may additionally enlarge a stabilizer in ambient dimensions, so effectiveness must be stated. Classifying all realizable subgroups and whether a group can arise through inequivalent actions or solutions remains open in this attempt. The first missing step is a collision-free dynamical existence criterion for each effective action, followed by exclusion of unwanted extra isotropy if full stabilizer realization is intended. The two interpretations in the source are kept separate.

\paragraph{Verification and adversarial audit.}
Checked pure-label collision argument, time/space cancellation, cyclic label action, mass preservation and ambient ineffective rotations. This was an LLM adversarial audit, not independent human verification; no theorem prover was used. Prompts A then B applied to this ID: Checked pure-label collision argument, time/space cancellation, cyclic label action, mass preservation and ambient ineffective rotations. Scope remains limited: Establish collision-free existence for prescribed effective actions and control extra isotropy when exact full stabilizers are requested.

Coordinator checked the collision-free pure-label kernel obstruction and combined time/rotation isotropy of a circle choreography. This additional check was model-assisted, not independent human verification.

\paragraph{Tools, process and contributions.}
Python 3 and local shell used to read supplied files and serialize this manuscript; mathematics derived symbolically in the text, with no numerical experiment or formal proof assistant. Contribution prompt A was followed by prompt B. GPT-6 Astra with ultra reasoning drafted the research and performed its stated self-audit; any additional coordinator contribution is identified in the verification and audit. A GPT-6-based Codex coordinator assembled the manuscript and checked repository integration; its exact variant and reasoning configuration were not exposed. No human mathematical contribution, independent human proof check or formal proof certification is supplied. Complete derivations and the consulted source references are included here; private reasoning transcripts are not published.

\paragraph{Final status and first remaining gap.}
\textbf{UNRESOLVED}. Establish collision-free existence for prescribed effective actions and control extra isotropy when exact full stabilizers are requested.

\section{Completion Estimate}
\noindent\textbf{COMPLETION: 0\%}
The calculation is an elementary benchmark or conditional reduction; it establishes no claimed new portion of the complete source target.

\subsection{Sources}
\begin{itemize}
\item[\hypertarget{source:20002542:1}{S1}] ULAM.AI and the UnsolvedMath Contributors, UnsolvedMath, record 20002542 (AIM-PHYSICS-0050). Catalogue curation under CC BY 4.0; definitions and English exposition added locally.
\url{https://huggingface.co/datasets/ulamai/UnsolvedMath}.
\item[\hypertarget{source:20002542:2}{S2}] American Institute of Mathematics, Variational Methods in Celestial Mechanics; notes by Kuo-Chang Chen, version 22 June 2003, source item 13.
\url{https://aimath.org/WWN/varcelest/varcelest.pdf}.
\end{itemize}
\label{end:20002542}
\end{document}
